\documentclass[11pt]{article}
\usepackage[a4paper,margin=1in]{geometry}
\usepackage{amsmath,amssymb,amsthm}
\usepackage{booktabs}
\usepackage{graphicx}
\usepackage{xcolor}
\usepackage{hyperref}
\hypersetup{colorlinks=true,linkcolor=blue,urlcolor=blue,citecolor=blue}
\newcommand{\rt}[1]{\texttt{#1}}
\newcommand{\fn}[1]{\texttt{#1}}
\newcommand{\Tr}{\operatorname{Tr}}
\newcommand{\Phid}{\Phi^{\dagger}}
\newcommand{\C}{\mathcal{C}}
\newcommand{\F}{\mathcal{F}}
\newcommand{\D}{\mathcal{D}}
\newcommand{\A}{\mathcal{A}}

\theoremstyle{plain}
\newtheorem{proposition}{Proposition}
\newtheorem{theorem}{Theorem}
\theoremstyle{remark}
\newtheorem{remark}{Remark}

\title{\textbf{Report R24 --- what the eight rules of R23 protect}\\
\large The strong-symmetry algebra, the conserved quantities, and why none of
it is a code}
\author{Tier 1 analysis}
\date{2026-08-22}

\begin{document}
\maketitle

\begin{abstract}
R23 established that the six bounded dissipative rules relax onto a wall tiling
with one period-2 dark qubit per wide tile, that $W_{29}/W_{71}$ split a
strong-symmetry block by a fair coin, and that $W_{73}/W_{109}$ instead reach a
hard-core constrained space carrying a Floquet blockade-Hadamard --- PXP with
the spin flip replaced by $H$. This report computes, exactly from the $4^N$
superoperator, the two algebras those statements live in: the strong-symmetry
commutant $\C$ and the conserved-quantity space $\F$. Three results. \emph{One},
the asymptotic algebra of every one of the eight is a \emph{single} full matrix
algebra $M_D$ with trivial mixing factor, $\dim\A = D^2$ exactly, where $D$ is
the total number of recurrent basis states; exhaustively at $N=4,5$, precisely
one Kraus operator is non-zero on the recurrent space and it is an isometry
there, so the dissipation is transient-only and the asymptotic evolution is a
closed unitary --- not merely one decoherence-free subspace per attractor, but a
single coherent block across all of them, with the coherences between distinct
attractors surviving too. \emph{Two}, R22's identity
$\dim(\F\cap\D) = n_{\mathrm{rec}}$ is a statement about the \emph{monitored}
chain; coherently only $n_{\mathrm{wcc}} \le \dim(\F\cap\D) \le
n_{\mathrm{rec}}$ survives, and $W_{29}$ at $N=4$ realises the strict case ---
its fair coin is not a conserved quantity. Auditing R22's eight gap rules the
same way splits them: five owe their gap entirely to the dephasing, while
$W_{36}, W_{44}, W_{100}$ keep a genuine coherent gap at $N=5$. \emph{Three},
none of the eight protects anything. Their code spaces are spanned by
computational basis states, and we prove that such a space admits no subspace at
all --- of any dimension above one --- that can so much as \emph{detect} a
single-site $Z$ error. Distance one, no subcode helps, verified against the
Knill--Laflamme conditions. What the eight are is a passive memory against one
specific structured noise model, namely the machine's own resets, which
annihilate the code space identically. The physics worth keeping is
$W_{73}/W_{109}$: dissipation autonomously prepares a Rydberg-blockade Hilbert
space of dimension $F(N+2)$ and then runs an entangling constrained Floquet
circuit inside it, with preparation efficiency decaying like $0.9^N$. That last
statement has a one-line origin, which we prove: $W_{73}$ and $W_{109}$ are the
unitary Floquet-$P^0HP^0$ and Floquet-$P^1HP^1$ rules $W_{201}$ and $W_{108}$
with a \emph{single identity gate replaced by a reset}, and on the constrained
space the reset can never fire, so the dissipative rule acts there exactly as
its unitary parent --- identically, to zero, for $N = 4,\dots,11$. This
generalises: all eight are children of the four single-$V$ rules, every
non-trivial enclosure is exactly a Krylov sector of its parent, and whether the
parent's projector reads a symmetric pattern or a domain wall decides the rest
--- symmetric gives the blockade space and an entangling PXP-type circuit,
asymmetric gives tiling subcubes carrying a bare tensor product of single-site
Hadamards. One slot of the gate word separates an interacting model from a free
one.
\end{abstract}

\tableofcontents

\section{The two algebras, and why they differ}

Throughout, $\{K_j\}$ are the Kraus operators of one brick-wall cycle, $\D$ is
the algebra of operators diagonal in the computational basis, and
\[
  \C = \{A : [A, K_j] = 0 \;\;\forall j\},
  \qquad
  \F = \{A : \Phid(A) = A\}.
\]
$\C$ is the strong-symmetry algebra of Bu\v{c}a and Prosen~\cite{buca2012}; $\F$
is the space of conserved quantities, the operational object, since
$\Tr(A\rho_t)$ is constant for every initial $\rho$ exactly when $A \in \F$.
Always $\C \subseteq \F$ --- verified here to $3\times10^{-14}$ --- and $\C$ is
an algebra while $\F$ in general is not: $\Phid$ respects products only when it
is a $*$-automorphism, i.e.\ only in the closed case. Frigerio's
theorem~\cite{frigerio1978} is the criterion: with a \emph{faithful} stationary
state $\F$ is a von Neumann algebra and equals $\C$. A large transient region is
precisely what denies that hypothesis, which is why the eight rules below have
$\dim\C$ far below $\dim\F$.

\section{The measurement}

\begin{table}[htbp]
\centering
\caption{The eight rules of R23, obc0, exact from the $4^N$ superoperator. $D$
is the dimension of the whole recurrent space, $\A$ the asymptotic algebra,
$\dim\mathrm{Fix}$ the fixed points of $\Phid$ inside it. Every row has
$\dim\A = D^2$.}
\label{tab:algebras}
\begin{tabular}{rlrrrrrrrr}
\toprule
rule & word & $N$ & $\dim\C$ & $\dim\F$ & $n_{\mathrm{wcc}}$ &
$n_{\mathrm{rec}}$ & $D$ & $\dim\A$ & $\dim\mathrm{Fix}$\\
\midrule
$W_{28}$  & IIVD & 4 & 25 & 45 & 6 & 6 &  9 &  81 &  45 \\
$W_{70}$  & IVID & 4 & 32 & 45 & 6 & 6 &  9 &  81 &  45 \\
$W_{29}$  & EIVD & 4 &  3 & 13 & 2 & \textbf{3} &  5 &  25 &  13 \\
$W_{71}$  & EVID & 4 &  6 & 13 & 3 & 3 &  5 &  25 &  13 \\
$W_{157}$ & EIVI & 4 &  3 & 13 & 3 & 3 &  5 &  25 &  13 \\
$W_{199}$ & EVII & 4 &  5 & 13 & 3 & 3 &  5 &  25 &  13 \\
$W_{73}$  & VIID & 4 & 19 & 35 & 4 & 4 & 13 & 169 &  35 \\
$W_{109}$ & EIIV & 4 &  9 & 31 & 5 & 5 &  9 &  81 &  31 \\
\midrule
$W_{28}$  & IIVD & 5 & 58 & 106 & 8 & 8 & 14 & 196 & 106 \\
$W_{70}$  & IVID & 5 & 58 & 106 & 8 & 8 & 14 & 196 & 106 \\
$W_{29}$  & EIVD & 5 & 11 &  41 & 4 & 4 &  9 &  81 &  41 \\
$W_{71}$  & EVID & 5 & 11 &  41 & 4 & 4 &  9 &  81 &  41 \\
$W_{157}$ & EIVI & 5 & 18 &  41 & 4 & 4 &  9 &  81 &  41 \\
$W_{199}$ & EVII & 5 & 18 &  41 & 4 & 4 &  9 &  81 &  41 \\
$W_{73}$  & VIID & 5 & 33 &  96 & 6 & 6 & 24 & 576 &  96 \\
$W_{109}$ & EIIV & 5 & 27 &  95 & 8 & 8 & 17 & 289 &  95 \\
\bottomrule
\end{tabular}
\end{table}

\subsection{The dissipation is transient-only}

\begin{proposition}[Exhaustive at $N = 4, 5$]
\label{prop:unitary}
For each of the eight rules, of the up to twelve Kraus operators of one cycle
exactly \emph{one} is non-zero on the recurrent subspace, and its restriction
$A$ there satisfies $\lVert A^\dagger A - \mathbb{1}\rVert_{\max} \le
3.6\times10^{-15}$. Every other Kraus operator annihilates that subspace
identically ($\lVert W^\dagger K W\rVert \le 7.6\times10^{-15}$).
\end{proposition}

\noindent Consequently the channel restricted to the recurrent space is a
unitary, and the asymptotic algebra is the full $M_D$ --- which is what
Table~\ref{tab:algebras} reports as $\dim\A = D^2$, in every row.

This is stronger than ``every attractor is a decoherence-free subspace''. That
statement is per-enclosure; this one is joint. The recurrent space of $W_{28}$
at $N=4$ carries six terminal components of sizes $2,2,2,1,1,1$, and the
channel decoheres neither within them \emph{nor between them}: a superposition
of states drawn from two different attractors is preserved exactly. It is also
what makes $\F$ readable, because with a unitary $U$ on the code space
\[
  \dim\F \;=\; \sum_i \mathrm{mult}(\lambda_i)^2 ,
\]
the sum over the multiplicities of $U$'s eigenvalues. Table~\ref{tab:algebras}
can then be read spectroscopically, and R23 rev.2 has since diagonalised $U$
directly, which both confirms the reading and corrects one guess made from the
totals alone.

For the six bounded rules $\dim\F \approx D^2/2$ --- $45/81$, $13/25$,
$106/196$, $41/81$ --- the signature of a two-fold spectrum. Diagonalising
confirms it exactly: two distinct eigenvalues every time and $U^2 = \mathbb{1}$
to $2.2\times10^{-16}$, with multiplicities $(6,3)$ for $W_{28}$ at $N=4$ and
$(3,2)$ for $W_{29}/W_{71}/W_{157}/W_{199}$ there, reproducing $45$ and $13$.
So R23's period-2 flicker, originally measured on trajectories, is forced by the
spectrum.

For $W_{73}$ at $N=5$ the total $\dim\F = 96$ is right but its factorisation is
\emph{not} $6\times4^2$, as a first reading of the totals suggested. $U$ has
twelve distinct eigenvalues with multiplicities $(8,4,2,2,1,1,1,1,1,1,1,1)$,
giving $64+16+4+4+8 = 96$; eight sit at exactly $+1$, four at $-1$, and the rest
are complex and off the roots of unity. $W_{109}$ at $N=5$ has $D=17$ and
multiplicities $(9,2,2,2,1,1)$, giving $95$. The conclusion is unchanged and in
fact sharpened --- complex eigenvalues off the unit roots mean no period at all,
$\lVert U^2 - \mathbb{1}\rVert = 0.5$ --- but a multiplicity pattern must be
read off the spectrum, never inferred from $\dim\F$.

\subsection{Scope}

Proposition~\ref{prop:unitary} covers all enclosures at $N = 4$ and $5$, which
is exhaustive at those sizes but only those sizes. The independent check on the
Tier-2 side (top three attractors per rule at $N=11$, plus a basis-by-basis pass
on the largest attractor of $W_{73}/W_{109}$ at $N=9,11$) reaches further up but
not across. Neither should be quoted as a theorem for all $N$.

\section{The squeeze: what the monitoring was doing}%
\label{sec:squeeze}

R22 Proposition~4 states $\dim(\F\cap\D) = n_{\mathrm{rec}}$. That is a
statement about the \emph{monitored} channel. For diagonal $A$ the diagonal of
$\Phid(A)$ is exactly the population chain's answer, so harmonicity is
necessary; but $\Phid(A)$ need not \emph{be} diagonal, and membership of
$\F\cap\D$ demands the off-diagonal part vanish as well. Hence
\begin{equation}
  n_{\mathrm{wcc}} \;=\; \dim(\C\cap\D)
  \;\le\; \dim(\F\cap\D) \;\le\; n_{\mathrm{rec}} .
  \label{eq:squeeze}
\end{equation}
The left equality is exact for the coherent channel and holds in every row we
have measured. The right inequality is saturated by construction for the
monitored chain and can be strict otherwise.

\paragraph{\texorpdfstring{$W_{29}$ at $N=4$}{W29 at N=4} is the counterexample.} Three terminal components,
and $\dim(\F\cap\D) = 2$ --- sharply, the smallest relative singular value of
the off-diagonal block is $0.357$, not a tolerance question. R23 identifies the
mechanism: the third component is reached from a shared transient by a fair
coin, and a fair coin is not a conserved quantity. The coherent channel simply
has no observable that records which way it fell.

\paragraph{Auditing R22's eight gap rules.} The same computation applied to the
rules that motivated R22 splits them in two (Table~\ref{tab:audit}). Five have
$\dim(\F\cap\D) = n_{\mathrm{wcc}}$ at both sizes: their WCC/SCC gap is an
artifact of dephasing and disappears when the monitoring is switched off. Three
--- $W_{36}, W_{44}, W_{100}$ --- keep a genuine coherent gap at $N=5$.

\begin{table}[htbp]
\centering
\caption{R22's gap rules re-measured coherently, obc0. ``monitored-only'' means
$\dim(\F\cap\D)$ collapses to $n_{\mathrm{wcc}}$ although the graph reports
$n_{\mathrm{rec}} > n_{\mathrm{wcc}}$.}
\label{tab:audit}
\begin{tabular}{rlrrrrl}
\toprule
rule & word & $N$ & $n_{\mathrm{wcc}}$ & $\dim(\F\cap\D)$ & $n_{\mathrm{rec}}$
& verdict \\
\midrule
$W_{44}$  & IIDV & 4 & 5 & 5 & 6 & monitored-only \\
$W_{104}$ & DIIV & 4 & 3 & 3 & 4 & monitored-only \\
$W_{203}$ & VEII & 4 & 1 & 1 & 3 & monitored-only \\
$W_{217}$ & VIEI & 4 & 1 & 1 & 3 & monitored-only \\
$W_{233}$ & VIIE & 4 & 3 & 3 & 4 & monitored-only \\
\midrule
$W_{104}$ & DIIV & 5 & 3 & 3 & 5 & monitored-only \\
$W_{203}$ & VEII & 5 & 1 & 1 & 4 & monitored-only \\
$W_{217}$ & VIEI & 5 & 1 & 1 & 4 & monitored-only \\
$W_{219}$ & VEEI & 5 & 1 & 1 & 2 & monitored-only \\
$W_{233}$ & VIIE & 5 & 3 & 3 & 4 & monitored-only \\
\midrule
$W_{36}$  & IDDV & 5 & 7 & \textbf{9} & 9 & genuine gap \\
$W_{44}$  & IIDV & 5 & 7 & \textbf{8} & 9 & genuine gap \\
$W_{100}$ & IDIV & 5 & 7 & \textbf{8} & 9 & genuine gap \\
\bottomrule
\end{tabular}
\end{table}

Two things this does \emph{not} disturb. R22's Theorem~1 concerns projections,
and $\C$ and $\F$ contain the same projections, so the Boolean algebra of
conserved questions still has exactly $n_{\mathrm{wcc}}$ atoms whichever channel
is meant. And the measurements of R22 \S3.4--3.5 are measurements of the
monitored chain, which is the right object for the question they answer --- what
a repeatedly observed system remembers. The correction is one of scope, not of
arithmetic.

\section{Trivial, and not}

\subsection{The six are physical qubits with extra steps}

$W_{28}, W_{70}, W_{29}, W_{71}, W_{157}, W_{199}$: R23 shows every enclosure is
a subcube $2^k$ over a frozen wall backbone, so the decoherence-free space
factorises as a tensor product of \emph{single-site} dark qubits, each running
$H^t$. There is no encoding: a logical qubit is one physical site. There is no
non-locality, hence nothing for a distance to be built out of, and the operator
entanglement of the asymptotic unitary is bounded, which is R23's mechanism for
the $N$-independent $\chi_{\mathrm{sat}}$.

$\dim\C$ for these is small and erratic --- $3, 5, 6, 25, 32$ at $N=4$, and note
that $W_{157}$ and $W_{199}$ have identical $\F$ but $\dim\C = 3$ against $5$.
That asymmetry is a fact about the \emph{transients}, since $\C$ must commute
with the Kraus operators everywhere, not only on the recurrent space. It is a
reminder that $\C$ is not the right invariant for asymptotic questions.

\subsection{\texorpdfstring{$W_{73}$ and $W_{109}$}{W73 and W109} are the interesting ones}

Their largest enclosure is the hard-core constrained space, of dimension
$F(N+2)$ and $F(N)$, and the induced one-period map is R23's blockade-controlled
Hadamard --- PXP with the spin flip replaced by $H$, in Floquet rather than
Hamiltonian form. Combining that with Proposition~\ref{prop:unitary}:

\begin{quote}
\textbf{The dissipation autonomously prepares a Rydberg-blockade Hilbert space
and then runs an entangling constrained Floquet circuit inside it.}
\end{quote}

That is a statement about dissipative state preparation of a kinetically
constrained manifold, and it is the physics in this family worth reporting. Its
efficiency is finite and decaying. The probability that a uniformly random
initial state is absorbed into the largest enclosure, from the exact absorption
probabilities:

\begin{center}
\begin{tabular}{lccccccccc}
\toprule
$N$ & 4 & 5 & 6 & 7 & 8 & 9 & 10 & 11 & 12 \\
\midrule
$W_{73}$  & $0.625$ & $0.563$ & $0.500$ & $0.438$ & $0.398$ & $0.344$ &
  $0.316$ & $0.272$ & $0.251$ \\
$W_{109}$ & $0.313$ & $0.438$ & $0.281$ & $0.328$ & $0.234$ & $0.254$ &
  $0.190$ & $0.199$ & $0.151$ \\
\midrule
$\dim$ & 8 & 13 & 21 & 34 & 55 & 89 & 144 & 233 & 377 \\
\bottomrule
\end{tabular}
\end{center}

\noindent Roughly $0.9^N$, against a target growing like $\varphi^N$. So the
preparation is probabilistic and needs heralding; it is not a deterministic
pump. Whether a modified rule can push this to $1$ is the obvious question and
is not attempted here.

\section{Nothing here is a code}
\label{sec:code}

\begin{theorem}[A basis-spanned code space has distance one]
\label{thm:distance}
Let the code space be spanned by computational basis states. Then no subspace of
it of dimension $\ge 2$ satisfies the Knill--Laflamme condition~\cite{kl1997}
for the single-site Pauli errors --- not the correction condition, and not even
the weaker detection condition.
\end{theorem}

\begin{proof}
Let $W$ be the isometry onto the code space and $V$ that of a candidate
subcode. Detection of $Z_i$ requires $V^\dagger W^\dagger Z_i W V = c_i
\mathbb{1}$, so the subcode lies inside an eigenspace of $M_i := W^\dagger Z_i
W$. Because $W$ is basis-spanned, each $M_i$ is diagonal in the codeword basis
with entries $\pm 1$, namely the $i$-th bit of each codeword. Distinct codewords
are distinct bit strings, so the joint spectrum $\{(\pm1)_i\}$ separates them
and every joint eigenspace of $M_1,\dots,M_N$ is one-dimensional. Hence
$\dim V \le 1$.
\end{proof}

R23 establishes that the asymptotic support of all eight rules is basis-spanned;
we re-verify it at $N = 4, 5$ against the superoperator. So all eight have
distance one and \emph{no subcode helps}. Checked directly: the Knill--Laflamme
condition fails at weight one for every one of them.

What is true is the complementary statement, and it is the same lesson R22 \S8
drew from majority voting. The eight are perfectly protected against the
machine's own dissipation --- Proposition~\ref{prop:unitary} says the jump
operators annihilate the code space identically, an \emph{infinite}-distance
noiseless subsystem in the sense of Knill, Laflamme and Viola~\cite{klv2000} for
that one structured noise model --- and completely unprotected against anything
else. Fragmentation protects everything it finds, including the errors, which is
why it protects nothing selectively.

One nuance worth recording, because it is the natural place to look for a
positive result. The blockade constraint of $W_{73}$ \emph{does} make some
bit-flip errors leave the code space, and those are detectable by a
constraint check. But not all of them: $\ldots0\,0\,0\ldots \to
\ldots0\,1\,0\ldots$ stays inside. The classical minimum distance of the
hard-core space is $1$, so even restricted to $X$ errors the constraint is a
detector of some errors, not a code.

\subsection{Three exponents, and no gap anywhere}
\label{sec:exponents}

Beyond the sizes the superoperator reaches, the graph engine gives the growth of
the protected space directly. It is worth being explicit about why that is
allowed: $D$ is computed by summing terminal-SCC sizes in the \emph{classical}
one-cycle support graph, and it equals the dimension of the decoherence-free
space only because Prop.~\ref{prop:unitary} says the whole recurrent space is
closed and carries a unitary. Without that, $D$ would be a combinatorial count
with no quantum meaning. With it, the graph reaches $N=18$ and each rule carries
\emph{three} distinct exponents:

\begin{table}[htbp]
\centering
\caption{Growth of the three quantities, obc0, from exact integer recurrences
verified term by term over $N = 4,\dots,18$. In both families the protected
dimension $D$ has the largest exponent --- the whole space grows faster than any
single enclosure in it.}
\label{tab:exponents}
\begin{tabular}{llll}
\toprule
 & enclosures $n_{\mathrm{rec}}$ & largest $d^{\mathrm{rec}}_{\max}$ &
   protected $D$ \\
\midrule
$W_{73}, W_{109}$ & $\psi = 1.4656$ & $\varphi = 1.6180$ &
  $\mathbf{1.8393}$, $x^3 = x^2+x+1$ \\
the six & $\rho = 1.3247$ & $2^{1/3}$ &
  $\mathbf{1.5214}$, $x^3 = x+2$ \\
\bottomrule
\end{tabular}
\end{table}

\noindent $D$ for $W_{73}$ runs $13, 24, 44, 81, 149, 274, 504, 927, 1705,
3136, \dots$, each term the sum of the previous three, ratio settled at
$1.8393$ by $N=9$; $W_{109}$ gives $9, 17, 31, 57, 105, 193, 355, \dots$, the
same recurrence. Among the six, $W_{28}/W_{70}$ give $9, 14, 20, 33, 49, 74,
116, \dots$ obeying $D_n = D_{n-1}+D_{n-2}+D_{n-3}-2D_{n-4}$, and
$W_{157}/W_{199}/W_{29}/W_{71}$ give $5, 9, 11, 19, 29, 41, 67, \dots$ obeying
$D_n = D_{n-2}+2D_{n-3}$ --- two different recurrences sharing the real root of
$x^3 = x+2$, a constant that appears nowhere else in this project. (The
recurrences and the identification of the shared root are the Tier-2 session's;
we verify all four term by term on independently generated series out to
$N=18$.)

Two readings. First, in both families the \emph{total} decoherence-free
dimension grows strictly faster than any single enclosure --- $0.879$ qubits per
site against $\log_2\varphi = 0.694$ for $W_{73}$, and $0.605$ against
$\tfrac13$ for the six. The excess is exactly the cross-enclosure coherence
Prop.~\ref{prop:unitary} guarantees, which a per-attractor account would miss
entirely. Second, all four of $W_{157}, W_{199}, W_{29}, W_{71}$ carry the
\emph{same} $D$ series although $W_{29}$ and $W_{71}$ have strictly fewer weak
components than enclosures: the multistability changes how the recurrent space
is partitioned into blocks, not how much of it there is.

One more column worth having: $n_{\mathrm{wcc}} = n_{\mathrm{rec}}$ at
\emph{every} $N$ from $4$ to $18$ for $W_{73}$ and $W_{109}$. Those two are
fragmented in both senses with no gap at any size we can reach --- which is what
a paper built on ``fragmented in both senses, and non-unitary'' needs, and which
\S\ref{sec:squeeze} shows is not automatic.

\section{\texorpdfstring{The dissipative pair is the unitary pair,\\ with one
letter changed}{The dissipative pair is the unitary pair, with one letter
changed}}
\label{sec:parents}

R23 identifies the induced unitary on the largest attractor of $W_{73}$ and
$W_{109}$ numerically, as a brickwork of blockade-controlled Hadamards. There is
a reason, and it is one line of gate arithmetic. It also holds for the other
six, which turns the two families into one statement (\S\ref{sec:allparents}).

Read the gate words in slot order $(00),(01),(10),(11)$, the slot being the
state of the two neighbours:
\[
  W_{201} = \mathrm{VIII}, \quad W_{73} = \mathrm{VIID};
  \qquad
  W_{108} = \mathrm{IIIV}, \quad W_{109} = \mathrm{EIIV}.
\]
$W_{201}$ fires a Hadamard exactly when both neighbours are empty --- it
\emph{is} the Floquet $P^0HP^0$ model --- and $W_{108}$ fires when both are
occupied, the $P^1HP^1$ model. And each dissipative rule is its unitary parent
with a \emph{single} identity replaced by a reset: $W_{73} = W_{201}$ with
slot $(11)$ changed $\mathrm{I}\to\mathrm{D}$, and $W_{109} = W_{108}$ with slot
$(00)$ changed $\mathrm{I}\to\mathrm{E}$.

\begin{proposition}[The reset is invisible on the constrained space]
\label{prop:parents}
Let $\mathcal{C}^1_N$ be the span of the basis states with no two adjacent
$1$s, and $\mathcal{C}^0_N$ likewise for $0$s. Then $\mathcal{C}^1_N$ is invariant
under the channel of $W_{73}$, no jump operator ever fires on it, and $W_{73}$
acts there exactly as $W_{201}$. Dually, $\mathcal{C}^0_N$ is invariant under
$W_{109}$, which acts there exactly as $W_{108}$.
\end{proposition}

\begin{proof}
$W_{73}$'s reset $\mathrm{D}$ sits in slot $(11)$, so it can only fire at a site
both of whose neighbours are $1$; on $\mathcal{C}^1_N$ such a site must itself
be $0$, and $\mathrm{D}$ resets to $0$, so it acts as the identity and its jump
Kraus operator vanishes. The only non-trivial gate left is the Hadamard in slot
$(00)$, which can set a site to $1$ only when both neighbours are $0$, creating
no adjacent pair; sites updated in the same brick-wall layer are two apart and
their common neighbour is $0$, so simultaneity does not spoil this. Hence
$\mathcal{C}^1_N$ is invariant and the map is the no-jump Kraus operator alone,
which is $W_{201}$'s unitary. The argument for $W_{109}$ is the same with $0$
and $1$, and $\mathrm{E}$ in place of $\mathrm{D}$.
\end{proof}

\noindent Measured, the two one-cycle matrices agree \emph{identically} --- not
to machine precision, to exactly zero --- for $N = 4,\dots,11$, with no
amplitude leaving the subspace, only the empty Kraus label occurring, and
$\lVert M^{\mathsf T}M - \mathbb{1}\rVert \le 2.3\times10^{-16}$.

So the mechanism behind ``some dissipative models inherit fragmentation'' is
not an analogy. The reset does no work at all in the steady state; its entire
role is to \emph{drive the system into} the constrained subspace, where the
dynamics is the unitary model itself. That also explains R23's $a=0$ versus
$a=1$ blockade: the projector is inherited from whichever slot the parent's
Hadamard occupies.

The proposition gives invariance, not attraction --- that $\mathcal{C}^1_N$ is
reached is a separate fact, and \S\ref{sec:squeeze}'s table shows it is reached
only with probability $\approx 0.9^N$. Irreducibility is a separate fact too,
and the two rules differ on it. Decomposing the constrained space (support
graph of the one-cycle map restricted to it, obc0):
\[
  \begin{array}{lll}
  W_{73} \text{ on no-}11: & N=9:\; \{89\} & N=11:\; \{233\}\\
  W_{109}\text{ on no-}00: & N=9:\; \{34,21,21,13\} & N=11:\; \{89,55,55,34\}
  \end{array}
\]
with no transient states inside the space at any $N$, all parts Fibonacci and
summing to $F(N+2)$. So $\mathcal{C}^1_N$ \emph{is} $W_{73}$'s attractor,
irreducibly, while for $W_{109}$ it splits four ways, the largest part being the
$F(N)$ that R23 reports.

It is worth being explicit about \emph{why}, because the plausible explanation
is wrong. One might guess the split comes from the \rt{obc0} boundary letting
$W_{109}$'s reset fire at the chain ends. It does not: the reset never fires
anywhere on $\mathcal{C}^0_N$, edges included. At the left edge the outside
counts as $0$, so site $0$ sees the pair $(0, x_1)$; if $x_1 = 0$ then $x_0$
must be $1$ to avoid a $00$, and $\mathrm{E}$ resets to $1$, the identity there;
if $x_1 = 1$ the slot is $(01) = \mathrm{I}$. Measured, exactly one Kraus label
occurs at every state for $N = 4,\dots,11$. The split is a property of the
\emph{parent}: $W_{108}$ is already reducible on that space, and gives the
identical four-way decomposition at every $N$ we checked. (Independently
confirmed on the Tier-2 side; the correction is theirs.)

\subsection{All eight, and what the fourth gate slot decides}
\label{sec:allparents}

The Tier-2 session asked whether Prop.~\ref{prop:parents} generalises. It does,
and the general form is more informative than the special case, because it makes
the split between the two families a property of \emph{one slot of the gate
word}.

The four unitary parents are exactly the four single-$V$ rules, differing only in
which neighbour pattern gates the Hadamard:
\[
  \begin{array}{llll}
  W_{201} = \mathrm{VIII} & \text{both empty} & P^0HP^0 &
    \text{parent of } W_{73}\\
  W_{108} = \mathrm{IIIV} & \text{both occupied} & P^1HP^1 &
    \text{parent of } W_{109}\\
  W_{156} = \mathrm{IIVI} & \text{left occupied, right empty} & P^1HP^0 &
    \text{parent of } W_{28}, W_{29}, W_{157}\\
  W_{198} = \mathrm{IVII} & \text{left empty, right occupied} & P^0HP^1 &
    \text{parent of } W_{70}, W_{71}, W_{199}
  \end{array}
\]

\begin{proposition}[The parent identity is general]
\label{prop:allparents}
For each of the eight, every non-trivial terminal SCC is exactly a Krylov sector
of its coherent parent, and on it the child's one-cycle matrix equals the
parent's, no jump operator fires, and the map is unitary.
\end{proposition}

\noindent Verified here on \emph{every} non-trivial terminal SCC at $N = 9$ and
$11$ --- $21$ and $41$ of them for $W_{28}/W_{70}$, $11$ and $20$ for
$W_{157}/W_{199}/W_{29}/W_{71}$, $24$ and $55$ for $W_{73}$, $27$ and $68$ for
$W_{109}$ --- with agreement \emph{identically zero}, zero escape, a single Kraus
label throughout, and $\lVert M^{\mathsf T}M-\mathbb{1}\rVert \le
2.3\times10^{-16}$. This upgrades R17, which had the same statement at the level
of the support graph, to the level of the operator.

The reading is then clean. \textbf{Whether the projector reads a symmetric
pattern or a domain wall decides everything else.} Symmetric, $(00)$ or $(11)$:
the gate is a blockade, the invariant space is the hard-core space of dimension
$F(N+2)$, and the dynamics entangles. Asymmetric, $(10)$ or $(01)$: the gate is a
domain-wall detector, the invariant spaces are the tiling subcubes, and the
dynamics does not entangle at all. One slot in the gate word separates a PXP-type
interacting model from a free one.

Free in the strongest sense: on every terminal SCC of the six the induced
unitary is \emph{exactly} the tensor product of single-site Hadamards over the
free sites, to $1.1\times10^{-16}$, at every SCC we checked. R23's period $2$ and
its exactly-zero entanglement both follow with no further measurement, since
$H^2=\mathbb{1}$ site by site and a product gate on a product subspace cannot
correlate anything.

That gives \S\ref{sec:exponents}'s table a structural reading. $D$ counts the
whole protected space, $d^{\mathrm{rec}}_{\max}$ counts one parent sector, and
the gap between the two exponents is coherence between \emph{different sectors of
the same parent} --- for the six, between different tilings of the same chain.
Neither family's protected space is a code, for the reason proved in
\S\ref{sec:code}; but the six's is not even a correlated state space. It is a
stack of independent qubits, one per wide tile per tiling.

\section{What this leaves for the PRL}

Three things the above supplies for a paper built on ``fragmented, in both
senses, and non-unitary''.

\begin{enumerate}
\item \textbf{The mechanism sentence, and it is better than expected.} The
dissipative rules inherit fragmentation because their dissipation is
\emph{transient-only}: asymptotically the channel is a unitary on an
exponentially large constrained subspace (Prop.~\ref{prop:unitary}). For the two
rules that matter it is sharper still --- $W_{73}$ and $W_{109}$ are the
Floquet-$P^0HP^0$ and Floquet-$P^1HP^1$ models with a \emph{single identity
replaced by a reset}, and on the constrained space the reset provably never
fires (Prop.~\ref{prop:parents}). The dissipative models do not merely resemble
the unitary ones; they \emph{are} them, plus a pump that drives the system into
the constrained sector, and by Prop.~\ref{prop:allparents} the same holds for
all eight, every enclosure being a Krylov sector of one of the four single-$V$
rules. If the paper introduces Floquet-PHP and Floquet-PHQ as its two unitary
models, then Fig.~3's dissipative counterparts are the same words with one
letter changed, which is about as compact as a third figure can get. That is
also why an SCC census suffices there, and why deferring the WCC/SCC distinction
to a companion paper is defensible: for seven of the eight, $n_{\mathrm{wcc}} =
n_{\mathrm{rec}}$ at every size we can check coherently. A bonus framing is
available too: the four parents differ in a single gate slot, and whether that
slot reads a symmetric pattern or a domain wall is exactly what separates the
interacting family from the free one.
\item \textbf{One caveat that must be handled before the figure is drawn.} It is
$W_{29}$, and now $W_{36}/W_{44}/W_{100}$ from \S\ref{sec:squeeze}: the SCC
count is not always a coherent invariant. A census presented as
``fragmentation'' should either restrict to rules where \eqref{eq:squeeze} is
tight or say which quantity is plotted.
\item \textbf{A gap in the evidence, and a trap in closing it.} R23 reports
$\chi_{\mathrm{sat}}$ saturating \emph{in $N$} for the six and growing
exponentially for $W_{73}/W_{109}$. A headline about linear growth \emph{in
time} is a different axis, and the dissipative rules have not been measured on
it. The trap is that the obvious protocol is self-defeating:
Proposition~\ref{prop:unitary} says there is no dissipation \emph{inside} an
enclosure, so a $\chi(t)$ run started from a state already in the largest
enclosure measures the operator entanglement of a purely unitary constrained
circuit and cannot support any claim about the dissipative model. The
dissipative content lives entirely in the transient. Whatever protocol is
chosen must therefore start outside the enclosure and pay for the transient.
A second, practical obstacle: the superoperator route OOMs for $W_{73}$ at
$N=7$, $t=5$ (27\,GB, because the zip-up materialises an untruncated bond at
$16\chi \approx 12000$), so those two need a contraction that never forms that
bond. $\chi(t)$ for the six bounded rules is runnable today.
\end{enumerate}

\section*{Reproduce}
\addcontentsline{toc}{section}{Reproduce}

Everything is in \fn{quantum/asymptotic.py}, which builds the dense $4^N$
superoperator through \fn{quantum/peripheral.py}. The commutant is intersected
one Kraus operator at a time rather than as a single stack --- with a reset at
every site the stack has $2^N d^2$ rows and exhausts memory at $N=5$.

\begin{verbatim}
    python -m qca_fragmentation.quantum.asymptotic --rules 28,29,73,109 --N 4
    pytest qca_fragmentation/tests/test_asymptotic.py   # 147 tests
\end{verbatim}

\noindent The tests pin Table~\ref{tab:algebras} row by row, check
Proposition~\ref{prop:unitary} for all eight at $N=4,5$, confirm the code space
is basis-spanned, verify Theorem~\ref{thm:distance} through the
Knill--Laflamme condition, and pin both sides of \eqref{eq:squeeze} including
the $W_{29}$ counterexample and the $W_{36}$-versus-$W_{203}$ split.

\begin{thebibliography}{9}

\bibitem{buca2012} B.~Bu\v{c}a and T.~Prosen,
  \emph{A note on symmetry reductions of the Lindblad equation},
  New J.\ Phys.\ \textbf{14}, 073007 (2012).

\bibitem{frigerio1978} A.~Frigerio,
  \emph{Stationary states of quantum dynamical semigroups},
  Commun.\ Math.\ Phys.\ \textbf{63}, 269 (1978).

\bibitem{kl1997} E.~Knill and R.~Laflamme,
  \emph{Theory of quantum error-correcting codes},
  Phys.\ Rev.\ A \textbf{55}, 900 (1997).

\bibitem{klv2000} E.~Knill, R.~Laflamme and L.~Viola,
  \emph{Theory of quantum error correction for general noise},
  Phys.\ Rev.\ Lett.\ \textbf{84}, 2525 (2000).

\bibitem{bknpv2008} R.~Blume-Kohout, H.~K.~Ng, D.~Poulin and L.~Viola,
  \emph{Characterizing the structure of preserved information in quantum
  processes},
  Phys.\ Rev.\ Lett.\ \textbf{100}, 030501 (2008).

\bibitem{iadecola2026} T.~Iadecola,
  \emph{Symmetry fragmentation},
  Phys.\ Rev.\ Lett.\ \textbf{136}, 100401 (2026), arXiv:2510.06333.

\end{thebibliography}

\end{document}
